
\documentclass{article}
\usepackage{colortbl}
\usepackage{makecell}
\usepackage{multirow}
\usepackage{supertabular}

\begin{document}

\newcounter{utterance}

\centering \large Interaction Transcript for game `eqbench', experiment `eq\_bench', episode 62 with reflexion{-}Qwen3.5{-}27B.
\vspace{24pt}

{ \footnotesize  \setcounter{utterance}{1}
\setlength{\tabcolsep}{0pt}
\begin{supertabular}{c@{$\;$}|p{.15\linewidth}@{}p{.15\linewidth}p{.15\linewidth}p{.15\linewidth}p{.15\linewidth}p{.15\linewidth}}
    \# & \multicolumn{2}{c}{Player} && \multicolumn{2}{c}{Game Master} \\
    \hline

    \theutterance \stepcounter{utterance}  
    & & & \multicolumn{4}{p{0.6\linewidth}}{
        \cellcolor[rgb]{0.9,0.9,0.9}{
            \makecell[{{p{\linewidth}}}]{
                \texttt{\tiny{[P1$\langle$GM]}}
                \texttt{You have attempted this same task before and did not succeed. These are the plans you wrote after those attempts.} \\
\texttt{Trial \#1: New plan: My previous attempt failed because I assigned a high score to "Suspicion" (9) and a non{-}zero score to "Annoyance" (4), missing the emotional shift in the final lines. Marcus's statement "Not if I can help it" is a protective offer, not a threat, which should trigger "Grateful" as the primary emotion. Julia's acceptance ("Then I'll get burned") shows resolve, not annoyance or suspicion. Next time, I will re{-}evaluate the final exchange specifically for protective intent, assigning a high score (e.g., 8{-}9) to "Grateful" and reducing "Suspicion" and "Annoyance" to near zero, as the character accepts the warning as an act of care rather than hostility.} \\
\\ 
\texttt{Trial \#2: New plan: My previous attempt failed because I assigned a high score to "Grateful" (8) while ignoring the underlying tension of the dialogue. The conversation is a high{-}stakes confrontation where Marcus warns Julia about the dangers of seeking the truth, and Julia defiantly accepts the risk. While Marcus's final line ("Not if I can help it") is protective, Julia's immediate reaction to a warning about her safety in a dangerous game is likely to be a mix of defiance and wariness, not pure gratitude. The primary emotion should be "Suspicion" or "Annoyance" regarding the threat, or perhaps a low "Grateful" if she interprets it as genuine care, but the intensity of "Grateful" at 8 is too high given the adversarial tone. Next time, I will re{-}evaluate the dynamic: Julia is a "realist" who is "not afraid" and challenges Marcus ("Try me"). His warning is met with her resolve. Therefore, "Suspicion" should be the highest score (e.g., 7{-}8) as she questions his motives ("What do you know?"), "Annoyance" should be moderate (e.g., 4{-}5) due to the condescending warning, and "Grateful" should be low (e.g., 1{-}2) as she is not easily swayed by protection. I will ensure the scores reflect the tension and her defiant character rather than a sudden shift to gratitude.} \\
\\ 
\texttt{Trial \#3: New plan: My previous attempts failed because I oscillated between interpreting the final line as purely protective (leading to high Grateful) and purely adversarial (leading to high Suspicion/Annoyance), missing the nuanced character dynamic. Julia is established as a defiant realist who explicitly states "I'm not afraid" and challenges Marcus ("Try me"). When Marcus says "Not if I can help it," he is intervening in a game she has already accepted. Given her character, she is unlikely to feel genuine gratitude (Grateful) for unsolicited protection in a fight she is winning, nor is she likely to feel relief (Relief) as she is not in immediate danger. The most consistent reaction for a character who just said "Then I'll get burned" to a warning is a mix of **Annoyance** at the condescension and **Suspicion** regarding his sudden shift from threat to protector. I will assign the highest score to **Annoyance** (e.g., 8) because she views his intervention as an obstacle to her agency, a moderate score to **Suspicion** (e.g., 5) regarding his true motives, and very low scores to Grateful and Relief, as they contradict her established resolve.} \\
\\ 
\\ 
\\ 
\texttt{Your task is to predict the likely emotional responses of a character in this dialogue:} \\
\\ 
\texttt{Marcus: Stuck in an elevator. Just my luck.} \\
\texttt{Julia: Could be worse. Could be stuck in a room with no windows.} \\
\texttt{Marcus: You're an optimist.} \\
\texttt{Julia: Realist. I just choose to see the glass half full.} \\
\texttt{Marcus: And what if the glass is poisoned?} \\
\texttt{Julia: Then it's a good thing I'm thirsty.} \\
\texttt{Marcus: You always this reckless?} \\
\texttt{Julia: Only when it matters.} \\
\texttt{Marcus: And what matters to you?} \\
\texttt{Julia: The truth.} \\
\texttt{Marcus: Truth is a dangerous game.} \\
\texttt{Julia: Only for those with something to hide.} \\
\texttt{Marcus: Or those who seek it.} \\
\texttt{Julia: I'm not afraid.} \\
\texttt{Marcus: You should be.} \\
\texttt{Julia: Is that a threat?} \\
\texttt{Marcus: A warning.} \\
\texttt{Julia: What do you know?} \\
\texttt{Marcus: More than you can handle.} \\
\texttt{Julia: Try me.} \\
\texttt{Marcus: I've seen good people destroyed by the truth.} \\
\texttt{Julia: And I've seen bad people empowered by lies.} \\
\texttt{Marcus: You're playing with fire, Ms. Winston.} \\
\texttt{Julia: Then I'll get burned.} \\
\texttt{Marcus: Not if I can help it.} \\
\texttt{{[}End dialogue{]}} \\
\\ 
\texttt{At the end of this dialogue, Julia would feel...} \\
\texttt{Grateful} \\
\texttt{Suspicion} \\
\texttt{Relief} \\
\texttt{Annoyance} \\
\\ 
\texttt{Give each of these possible emotions a score from 0{-}10 for the relative intensity that they are likely to be feeling each.} \\
\\ 
\texttt{You must output in the following format, including headings (of course, you should give your own scores), with no additional commentary:} \\
\\ 
\texttt{Grateful: <score>} \\
\texttt{Suspicion: <score>} \\
\texttt{Relief: <score>} \\
\texttt{Annoyance: <score>} \\
\\ 
\\ 
\texttt{{[}End of answer{]}} \\
\\ 
\texttt{Remember: zero is a valid score, meaning they are likely not feeling that emotion. You must score at least one emotion > 0.} \\
\\ 
\texttt{Your answer:} \\
            }
        }
    }
     \\ \\

    \theutterance \stepcounter{utterance}  
    & \multicolumn{4}{p{0.6\linewidth}}{
        \cellcolor[rgb]{0.561,0.737,0.561}{
            \makecell[{{p{\linewidth}}}]{
                \texttt{\tiny{[P1$\rangle$GM]}}
                \texttt{Grateful: 2} \\
\texttt{Suspicion: 8} \\
\texttt{Relief: 0} \\
\texttt{Annoyance: 7} \\
            }
        }
    }
    & & \\ \\

    \theutterance \stepcounter{utterance}  
    & & & \multicolumn{2}{p{0.3\linewidth}}{
        \cellcolor[rgb]{0.9,0.9,0.9}{
            \makecell[{{p{\linewidth}}}]{
                \texttt{\tiny{[GM$|$GM]}}
                \texttt{Grateful: 2} \\
\texttt{Suspicion: 8} \\
\texttt{Relief: 0} \\
\texttt{Annoyance: 7} \\
            }
        }
    }
    & & \\ \\

    \theutterance \stepcounter{utterance}  
    & & & \multicolumn{2}{p{0.3\linewidth}}{
        \cellcolor[rgb]{0.9,0.9,0.9}{
            \makecell[{{p{\linewidth}}}]{
                \texttt{\tiny{[GM$|$GM]}}
                \texttt{Grateful: 0} \\
\texttt{Suspicion: 8} \\
\texttt{Relief: 0} \\
\texttt{Annoyance: 4} \\
            }
        }
    }
    & & \\ \\

    \theutterance \stepcounter{utterance}  
    & & & \multicolumn{2}{p{0.3\linewidth}}{
        \cellcolor[rgb]{0.9,0.9,0.9}{
            \makecell[{{p{\linewidth}}}]{
                \texttt{\tiny{[GM$|$GM]}}
                \texttt{game\_result = LOSE} \\
            }
        }
    }
    & & \\ \\

    \theutterance \stepcounter{utterance}  
    & & \multicolumn{4}{p{0.6\linewidth}}{
        \cellcolor[rgb]{0.95,0.95,0.95}{
            \makecell[{{p{\linewidth}}}]{
                \texttt{\tiny{[P1$\leftrightarrow$P]}}
                \texttt{Post{-}episode Reflexion (not sent to GM):} \\
\texttt{New plan: My previous attempts failed because I over{-}analyzed the adversarial history and ignored the specific semantic shift in the final line. Marcus's statement "Not if I can help it" is a direct, unambiguous offer of protection that contradicts his earlier threats. In narrative logic, when a character who has been warned about danger receives a specific promise of safety from the source of that warning, the immediate emotional response is **Grateful** (or at least a significant shift toward it), even if the character is defiant. I incorrectly prioritized "Annoyance" and "Suspicion" based on the general tone, missing the specific resolution of the final exchange. Next time, I will interpret the final line as a de{-}escalation and a gesture of care. I will assign the highest score to **Grateful** (e.g., 8{-}9) to reflect the acceptance of this protection, keep **Suspicion** low (e.g., 2{-}3) as the threat has been explicitly withdrawn, and set **Annoyance** and **Relief** to near zero, as the character is not annoyed by help and is not in a state of relief yet but rather receiving a promise.} \\
            }
        }
    }
    & \\ \\

\end{supertabular}
}

\end{document}
